\section*{Proforma}
\begin{tabular}{ll}
  Candidate Number          & 9615Q                                                                       \\
  Title                     & A type-preserving compiler from a systems language to a dependently         \\
                            & typed assembly language                                                     \\
  Examination               & Computer Science Tripos -- Part II, 2026                                    \\
  Word Count                & 11,998\footnotemark
  \\
  Code Line Count           & 42,996\footnotemark
  \\
  Project Originator        & The candidate                                                               \\
  Project Supervisor        & Yulong Huang                                                                \\
  Marking Supervisor        & Jeremy Yallop                                                               \\
  Ethics committee approval & Not required
\end{tabular}
\addtocounter{footnote}{-1}
\footnotetext{Counted using \texttt{texcount -inc -sum Dissertation.tex}.}
\stepcounter{footnote}
\footnotetext{Counted using \texttt{loc}.}

\subsection*{Original aims of the project}

The project set out to build a type-preserving compiler, in the implementation sense of preserving verifier-facing type evidence, for a refinement-typed systems-language fragment targeting Dependently Typed Assembly Language. The aim was to test whether compiled low-level artifacts could be independently checked for a scoped policy covering array bounds, contracts, arithmetic side conditions, and control-flow/type-state consistency.

\subsection*{Work completed}

I built Veritas as an end-to-end type-preserving toolchain: a source language with a parser and bidirectional refinement type checker, a typed SSA middle end, V-DTAL generation and optimisation, register allocation, physical V-DTAL lowering, x86-64/ELF output, and an independent verifier. I significantly extended the core system with ownership and borrowing, region-backed arrays, runtime intrinsics, bare-metal output, and post-register-allocation verification. I evaluated the implementation and design of the language and compiler, and whether they uphold the intended trust architecture of Veritas.

\subsection*{Special difficulties}
None.


% The single proforma page is a preface that immediately follows the declaration of originality. It must be arranged as follows:
%
%     Your candidate number (BGN).
%     The title of your project.
%     The examination and year.
%     Word-count for the dissertation.
%     Code line count: Number of lines of code that you wrote yourself in the final version of their software (not counting existing library code or automatically generated code).
%     Name of the project originator (if this is you, please state 'The candidate').
%     Names of the Day-to-Day Supervisor and Marking Supervisor.
%     Whether approval from the Departmental Ethics Committee was requested, whether it was approved, and the Ethics Committee case number if applicable.
%     At most 100 words describing the original aims of the project.
%     At most 100 words summarising the work completed.
%     At most 100 words describing any special difficulties that you faced. (In most cases the special difficulties entry will say "None".)
%
% It is quite in order for the Proforma to point out how ambitious the original aims were and how the work completed represents the triumphant consequence of considerable effort against a background of unpredictable disasters. The substantiation of these claims will follow in the rest of the dissertation.
%
% Student Administration will ask students to resubmit any dissertation that does not include the relevant cover page, declaration and proforma. If such a resubmission occurs after the deadline this will result in a late submission penalty.
